\documentclass[11pt,a4paper]{article}

% Pakete
\usepackage[utf8]{inputenc}
\usepackage[T1]{fontenc}
\usepackage[ngerman]{babel}
\usepackage[left=18mm,right=18mm,top=18mm,bottom=18mm]{geometry}
\usepackage{helvet}
\renewcommand{\familydefault}{\sfdefault}
\usepackage{amsmath,amssymb}
\usepackage{graphicx}
\usepackage{tikz}
\usetikzlibrary{arrows.meta,positioning}
\usepackage{tcolorbox}
\tcbuselibrary{skins,breakable}
\usepackage{enumitem}
\usepackage{tabularx}
\usepackage{colortbl}
\usepackage{qrcode}
\usepackage{xcolor}
\usepackage[hidelinks]{hyperref}

% Fachfarbe Physik
\definecolor{physik}{HTML}{2c5aa0}
\definecolor{physiklight}{HTML}{e8f0fa}

% Box-Stile
\tcbset{
    formelbox/.style={
        colback=white,
        colframe=physik,
        boxrule=2pt,
        arc=2pt,
        left=8pt,
        right=8pt,
        top=8pt,
        bottom=8pt
    }
}

% Kopfzeile
\usepackage{fancyhdr}
\pagestyle{fancy}
\fancyhf{}
\lhead{\small Physik 12 GK}
\chead{\small Bohr-Atommodell: Rechnen}
\rhead{\small Name: \underline{\hspace{4cm}}}
\renewcommand{\headrulewidth}{0.4pt}

% Keine Einrueckung
\setlength{\parindent}{0pt}
\setlength{\parskip}{6pt}

\begin{document}

% Titel
\begin{minipage}[t]{0.65\textwidth}
\vspace{0pt}
{\Large\textbf{Bohr-Atommodell: Rechnen}}\\[0.3em]
{\normalsize Fr 30.01.2026 | Physik 12 GK}
\end{minipage}
\hfill
\begin{minipage}[t]{0.32\textwidth}
\vspace{0pt}
\raggedleft
\begin{tabular}{@{}c@{}}
\href{https://devbox42.github.io/sciencesim/physik/kl12-GK/01-photoeffekt/06-ES-bohr-postulate/LP-06-bohr-postulate.html}{\qrcode[height=1.1cm]{https://devbox42.github.io/sciencesim/physik/kl12-GK/01-photoeffekt/06-ES-bohr-postulate/LP-06-bohr-postulate.html}} \\
{\scriptsize Lernpfad}
\end{tabular}
\end{minipage}

\vspace{0.3em}
\hrule height 1pt
\vspace{0.6em}

% ============================================================
% FORMELKASTEN
% ============================================================

\begin{tcolorbox}[formelbox]
\textbf{Formeln zum Bohr-Modell}

\vspace{0.5em}

\begin{minipage}[t]{0.48\textwidth}
\textbf{Energieniveaus:}
\begin{center}
{\large $E_n = \dfrac{-13{,}6 \text{ eV}}{n^2}$}
\end{center}
\end{minipage}
\hfill
\begin{minipage}[t]{0.48\textwidth}
\textbf{Rydberg-Formel:}
\begin{center}
{\large $\dfrac{1}{\lambda} = R_H \cdot \left( \dfrac{1}{n_1^2} - \dfrac{1}{n_2^2} \right)$}
\end{center}
\end{minipage}

\vspace{0.5em}

\textbf{Photonenenergie:} $E_{\text{Photon}} = |E_{n_1} - E_{n_2}| = h \cdot f$ \hfill
\textbf{Wellenlänge:} $\lambda = \dfrac{h \cdot c}{E}$ mit $h \cdot c = 1240 \text{ eV} \cdot \text{nm}$

\vspace{0.3em}

\textbf{Rydberg-Konstante:} $R_H = 1{,}097 \cdot 10^7 \text{ m}^{-1}$ \hfill
\textbf{Für wasserstoffähnliche Ionen:} $E_n = \dfrac{-13{,}6 \text{ eV} \cdot Z^2}{n^2}$
\end{tcolorbox}

\vspace{0.5em}

% ============================================================
% AUFGABE 1: Energieniveaus
% ============================================================

\textbf{Aufgabe 1: Energieniveaus berechnen} \hfill (3 BE)

Berechne die Energien $E_1$, $E_2$, $E_3$ und $E_4$ für das Wasserstoffatom.

\vspace{0.3em}

\begin{center}
\begin{tabular}{|c|c|c|c|c|c|}
\hline
$n$ & 1 & 2 & 3 & 4 & $\infty$ \\
\hline
$E_n$ (eV) & & & & & 0 \\
\hline
\end{tabular}
\end{center}

\vspace{1.5cm}

% ============================================================
% AUFGABE 2: Emission
% ============================================================

\textbf{Aufgabe 2: Emission -- Übergang $n = 3 \to n = 2$} \hfill (4 BE)

a) Berechne die Energie des emittierten Photons. (2 BE)

\vspace{1.8cm}

b) Berechne die Wellenlänge dieses Photons und gib die Farbe an. (2 BE)

\vspace{1.8cm}

% ============================================================
% AUFGABE 3: Rydberg-Formel
% ============================================================

\textbf{Aufgabe 3: Rydberg-Formel anwenden} \hfill (3 BE)

Berechne die Wellenlänge des Übergangs $n = 4 \to n = 2$ direkt mit der Rydberg-Formel.

\vspace{3cm}

% ============================================================
% AUFGABE 4: Absorption
% ============================================================

\textbf{Aufgabe 4: Absorption} \hfill (3 BE)

Ein Wasserstoffatom im Grundzustand ($n = 1$) absorbiert ein Photon und das Elektron springt auf $n = 3$.

a) Welche Energie muss das Photon haben? (2 BE)

\vspace{1.5cm}

b) Welche Wellenlänge hat dieses Photon? In welchem Bereich des Spektrums liegt es? (1 BE)

\vspace{1.5cm}

% ============================================================
% AUFGABE 5: Transfer He+ (AFB III)
% ============================================================

\textbf{Aufgabe 5: Transfer -- Helium-Ion He$^+$} \hfill (3 BE)

Das einfach ionisierte Helium (He$^+$) hat nur ein Elektron, aber $Z = 2$.

Berechne die Wellenlänge des Photons beim Übergang $n = 3 \to n = 2$ in He$^+$.

\vspace{3.5cm}

% ============================================================
% AUFGABE 6: Grenzen (AFB III)
% ============================================================

\textbf{Aufgabe 6: Grenzen des Bohr-Modells} \hfill (2 BE)

Das Bohr-Modell kann das Spektrum von Helium (He, zwei Elektronen) nicht korrekt vorhersagen.

Erkläre, warum das Modell hier versagt und beurteile, ob es trotzdem sinnvoll ist, das Modell im Unterricht zu verwenden.

\vspace{2.5cm}

% ============================================================
% AUFGABE 7: Spektralanalyse
% ============================================================

\textbf{Aufgabe 7: Spektralanalyse} \hfill (2 BE)

Öffne die Simulation zur Spektralanalyse im Lernpfad. Wähle den \textbf{Astronomie-Modus} und analysiere das \textbf{Sonnenspektrum}.

\vspace{0.5em}

a) Welche Elemente kannst du in der Sonnenatmosphäre nachweisen? (1 BE)

\vspace{1.5cm}

b) Erkläre, warum im Sonnenspektrum \emph{Absorptionslinien} (dunkel) statt Emissionslinien (hell) zu sehen sind. (1 BE)

\vspace{2cm}

\vfill

\begin{center}
\rule{10cm}{0.5pt}

\textbf{Gesamt:} \underline{\hspace{1cm}} / 20 BE
\end{center}

\end{document}
